\documentclass[11pt,a4paper]{article}
% ==========================================================
%  PREÁMBULO — Ecuaciones Diferenciales, Tema 2
%  Motor: pdfLaTeX · Fuente: Palatino (newpx), como el artículo de referencia
% ==========================================================
\usepackage[utf8]{inputenc}
\usepackage[T1]{fontenc}
\usepackage[spanish,es-tabla,es-nodecimaldot]{babel}
\usepackage[p,osf]{newpxtext}
\usepackage{amsmath, amsthm, mathtools}
\usepackage[bigdelims,vvarbb]{newpxmath}
\usepackage[final,protrusion=true,expansion=true]{microtype}
\usepackage{bm}
\usepackage{siunitx}
\usepackage{graphicx}
\usepackage{enumitem}
\usepackage{booktabs}
\usepackage{array}
\usepackage{float}
\usepackage{fancyhdr}
\usepackage{needspace}
\usepackage{titlesec}
\usepackage{xcolor}
\usepackage{caption}
\usepackage[a4paper,left=3cm,right=3cm,top=3cm,bottom=3cm]{geometry}

\sisetup{locale=FR, per-mode=symbol, inter-unit-product=\ensuremath{{}\cdot{}}}
\sisetup{output-decimal-marker={,}}

% --- Paleta sobria (azules y grises) ---
\definecolor{azulOscuro}{HTML}{1F3A5F}
\definecolor{azulMedio}{HTML}{3B6EA5}
\definecolor{azulClaro}{HTML}{8FB3D9}
\definecolor{azulFondo}{HTML}{F3F7FB}
\definecolor{gris}{HTML}{6B7280}
\definecolor{grisClaro}{HTML}{D4D8DE}
\definecolor{acento}{HTML}{B5651D}

% --- Enlaces e índice ---
\usepackage[colorlinks=true,linkcolor=azulOscuro,citecolor=azulOscuro,urlcolor=azulMedio]{hyperref}
\usepackage[nameinlink,noabbrev,spanish]{cleveref}
\usepackage{tocloft}
\renewcommand{\cfttoctitlefont}{\hfill\large\scshape\color{azulOscuro}}
\renewcommand{\cftaftertoctitle}{\hfill}
\renewcommand{\cftsecleader}{\hfill}
\renewcommand{\cftsubsecleader}{\hfill}
\renewcommand{\cftsecfont}{\color{azulMedio}}
\renewcommand{\cftsubsecfont}{\color{azulMedio}}
\renewcommand{\cftsecpagefont}{\color{azulMedio}}
\renewcommand{\cftsubsecpagefont}{\color{azulMedio}}
\setcounter{tocdepth}{2}

% --- Tipografía y página ---
\linespread{1.08}
\setlength{\parindent}{1.2em}
\setlength{\parskip}{0pt}
\setlist{itemsep=2pt, topsep=4pt}
\pagestyle{fancy}
\fancyhf{}
\renewcommand{\headrulewidth}{0pt}
\setlength{\headheight}{14pt}
\fancyhead[C]{\footnotesize\scshape\color{gris} Ecuaciones diferenciales · Tema 2}
\fancyhead[R]{\footnotesize\color{gris}\thepage}
\fancypagestyle{plain}{\fancyhf{}\fancyfoot[C]{\footnotesize\color{gris}\thepage}\renewcommand{\headrulewidth}{0pt}}

% Estilo AMS: sección centrada en versalitas; subsección en negrita
\titleformat{\section}[block]{\centering\scshape\color{azulOscuro}}{\thesection.}{0.5em}{}
\titleformat{\subsection}[block]{\bfseries\color{azulOscuro}}{\thesubsection.}{0.5em}{}
\titleformat{name=\subsection,numberless}[block]{\itshape\color{azulOscuro}}{}{0pt}{}
\titlespacing*{\section}{0pt}{3ex plus 1ex}{1.6ex}
\titlespacing*{\subsection}{0pt}{2.2ex plus .8ex}{1ex}

% --- Notación ---
\newcommand{\dd}{\mathrm{d}}
\newcommand{\ee}{\mathrm{e}}
\newcommand{\dv}[2]{\dfrac{\dd #1}{\dd #2}}
\newcommand{\R}{\mathbb{R}}
\newcommand{\abs}[1]{\lvert #1 \rvert}
\newcommand{\eqnum}[1]{\tag{#1}}
\mathtoolsset{showonlyrefs=false}
\numberwithin{equation}{section}

% --- Gráficas ---
\usepackage{tikz}
\usetikzlibrary{arrows.meta, calc, positioning, decorations.markings}
\usepackage[siunitx]{circuitikz}
\usepackage{pgfplots}
\pgfplotsset{compat=1.17}
\pgfplotsset{
  informe/.style={
    width=0.9\linewidth, height=6.4cm,
    axis lines=left, axis line style={gris},
    tick style={gris}, tick label style={font=\footnotesize},
    label style={font=\small},
    grid=major, grid style={grisClaro, very thin},
    enlargelimits=false, clip=true,
    legend style={font=\footnotesize, draw=none, fill=none, at={(0.97,0.5)}, anchor=east, legend cell align=left},
  },
  serie1/.style={color=azulOscuro, line width=1.1pt},
  serie2/.style={color=azulMedio, line width=1.1pt, dashed},
  serie3/.style={color=gris, line width=1.1pt, dotted},
  serie4/.style={color=acento, line width=1.1pt, dash dot},
  asintota/.style={color=gris, line width=0.6pt, dashed},
}
\captionsetup{font=small, labelfont={bf,color=azulOscuro}, labelsep=period}

% ==========================================================
%  ENTORNOS
% ==========================================================
\usepackage[most]{tcolorbox}

\newcommand{\etiquetaEntorno}[1]{\textsc{#1}}

% Teoremas, proposiciones, lemas, corolarios: cuerpo en cursiva, encabezado azul
\newtheoremstyle{enunciado}{9pt}{9pt}{\itshape}{}{\bfseries\color{azulOscuro}}{.}{0.6em}%
  {\thmname{#1}\ \thmnumber{#2}\thmnote{\ (#3)}}
% Ejemplos y observaciones: cuerpo recto, encabezado en cursiva azul
\newtheoremstyle{nota}{8pt}{8pt}{\normalfont}{}{\itshape\bfseries\color{azulMedio}}{.}{0.6em}%
  {\thmname{#1}\ \thmnumber{#2}\thmnote{\ (#3)}}

\theoremstyle{enunciado}
\newtheorem{theorem}{Teorema}[section]
\newtheorem{proposition}[theorem]{Proposición}
\newtheorem{lemma}[theorem]{Lema}
\newtheorem{corollary}[theorem]{Corolario}
\theoremstyle{nota}
\newtheorem{example}{Ejemplo}[section]
\newtheorem{remark}[example]{Observación}

\renewcommand{\qedsymbol}{\textcolor{azulMedio}{$\blacksquare$}}
\renewcommand{\proofname}{\itshape Demostración}
\newcommand{\finejemplo}{\hfill\textcolor{azulMedio}{$\diamond$}}

% Definiciones: único entorno enmarcado (fino, fondo tenue)
\newtcbtheorem[number within=section]{definition}{Definición}{
  enhanced, breakable, sharp corners,
  colback=azulFondo, colframe=azulClaro, boxrule=0.4pt,
  left=9pt, right=9pt, top=5pt, bottom=5pt,
  before skip=10pt, after skip=10pt,
  fonttitle=\bfseries\color{azulOscuro}, coltitle=azulOscuro,
  attach boxed title to top left={yshift*=-\tcboxedtitleheight/2, xshift=8pt},
  boxed title style={colback=white, colframe=white, sharp corners, boxrule=0pt, left=3pt, right=3pt, top=1pt, bottom=1pt},
  title={}}{def}
\newenvironment{definicion}[2][]{\begin{definition}{#2}{#1}}{\end{definition}}

% Modelo: filete superior e inferior
\newtcolorbox{modelo}[1][]{
  enhanced, breakable, blanker,
  borderline north={0.6pt}{0pt}{azulOscuro},
  borderline south={0.6pt}{0pt}{azulOscuro},
  top=6pt, bottom=6pt, left=2pt, right=2pt,
  before skip=10pt, after skip=10pt,
  before upper={\noindent{\scshape\bfseries\color{azulOscuro} Modelo\if\relax\detokenize{#1}\relax\else\ --- #1\fi.}\ }}

% Desarrollo teórico: línea lateral fina gris
\newtcolorbox{desarrollo}[1][]{
  enhanced, breakable, blanker,
  borderline west={0.5pt}{0pt}{grisClaro},
  left=12pt, top=2pt, bottom=2pt,
  before skip=8pt, after skip=8pt,
  before upper={\noindent{\scshape\color{gris} Desarrollo\if\relax\detokenize{#1}\relax\else\ --- #1\fi.}\ }}

% Algoritmo de resolución: pasos numerados
\newcounter{algoritmo}[section]
\renewcommand{\thealgoritmo}{\thesection.\arabic{algoritmo}}
\newtcolorbox{algoritmo}[1][]{
  enhanced, breakable, blanker,
  borderline north={0.6pt}{0pt}{azulOscuro},
  borderline south={0.3pt}{0pt}{azulClaro},
  top=6pt, bottom=6pt, left=2pt, right=2pt,
  before skip=10pt, after skip=10pt,
  before upper={\refstepcounter{algoritmo}\noindent{\scshape\bfseries\color{azulOscuro} Algoritmo \thealgoritmo\if\relax\detokenize{#1}\relax\else\ (#1)\fi.}\par\vspace{2pt}},
  }
\newlist{pasos}{enumerate}{1}
\setlist[pasos]{label=\textbf{\textcolor{azulMedio}{Paso \arabic*.}}, leftmargin=3.6em, labelsep=0.5em, itemsep=2pt}

% Ejercicios
\newcounter{ejercicio}[section]
\renewcommand{\theejercicio}{\thesection.\arabic{ejercicio}}
\newtcolorbox{ejercicio}[1][]{
  enhanced, breakable, blanker,
  borderline west={1.2pt}{0pt}{azulMedio},
  left=12pt, top=3pt, bottom=3pt,
  before skip=10pt, after skip=10pt,
  before upper={\refstepcounter{ejercicio}\noindent{\bfseries\color{azulOscuro} Ejercicio \theejercicio\if\relax\detokenize{#1}\relax\else\ (#1)\fi.}\ }}

% Tabla de parámetros
\newcommand{\encabezadoTabla}[1]{\textbf{\textcolor{azulOscuro}{#1}}}

\crefname{figure}{figura}{figuras}
\Crefname{figure}{Figura}{Figuras}
\crefname{equation}{ecuación}{ecuaciones}
\Crefname{equation}{Ecuación}{Ecuaciones}
\crefname{theorem}{teorema}{teoremas}
\crefname{proposition}{proposición}{proposiciones}
\crefname{lemma}{lema}{lemas}
\crefname{algoritmo}{algoritmo}{algoritmos}

% Resumen estilo AMS
\newenvironment{resumen}{\small\begin{quote}\noindent\textsc{Resumen.}\ }{\end{quote}}

% --- Utilidades adicionales ---
\usepackage{subcaption}
\captionsetup[subfigure]{font=footnotesize, labelfont={color=azulOscuro}, labelsep=period}
\newcommand{\metodo}[1]{\par\smallskip\noindent\textbf{\textcolor{azulMedio}{#1.}}\ }
\newenvironment{solucion}{\par\smallskip\noindent\textit{\textcolor{azulOscuro}{Solución.}}\ }{\hfill\textcolor{azulMedio}{$\diamond$}\par\smallskip}
\newcommand{\Tf}{t_{\mathrm{f}}}

% Problemas de seminario
\newcounter{problema}[section]
\renewcommand{\theproblema}{\thesection.\arabic{problema}}
\newtcolorbox{problema}[1][]{
  enhanced, breakable, blanker,
  borderline west={1.2pt}{0pt}{azulMedio},
  left=12pt, top=3pt, bottom=3pt,
  before skip=12pt, after skip=6pt,
  before upper={\refstepcounter{problema}\noindent{\bfseries\color{azulOscuro} Problema \theproblema\if\relax\detokenize{#1}\relax\else\ (#1)\fi.}\ }}


\begin{document}

% ==========================================================
%  CABECERA
% ==========================================================
\begin{center}
  \includegraphics[height=1.4cm]{figuras/logo_usal.png}\\[1.2em]
  {\color{azulOscuro}\large\scshape Aplicaciones de las ecuaciones diferenciales ordinarias\\ de primer orden a modelos sencillos}\\[1em]
  {\footnotesize\scshape Daniel Martín Guerrero}\\[0.2em]
  {\footnotesize Ecuaciones Diferenciales · Grado en Física · Facultad de Ciencias · Curso 2026--2027}
\end{center}
\begin{resumen}
Se estudian modelos de la física, la biología y la ingeniería que se reducen a ecuaciones diferenciales ordinarias de primer orden: enfriamiento de Newton, dinámica de poblaciones, competición de especies, fuerzas dependientes de la velocidad y circuitos eléctricos.
\end{resumen}
\vspace{0.5em}

\tableofcontents
\vspace{0.5em}

\section{Ley de Newton del enfriamiento}

La ley de Newton describe el intercambio de calor de un cuerpo con su entorno cuando domina la convección y la diferencia de temperaturas es moderada: la tasa de variación de la temperatura es proporcional a dicha diferencia.

\begin{modelo}[enfriamiento de Newton]
Sean $T(t)$ la temperatura del cuerpo y $S(t)$ la del entorno. Entonces
\begin{equation}
  \dv{T}{t} = -k\,\bigl(T - S\bigr), \qquad T(0)=T_0,
  \label{eq:newton}
\end{equation}
con $k>0$ constante de enfriamiento, $[k]=\si{\per\second}$.
\end{modelo}

\begin{remark}
El signo de \eqref{eq:newton} expresa que el calor fluye del cuerpo más caliente al más frío: si $T>S$ entonces $\dot T<0$, y si $T<S$ entonces $\dot T>0$.
\end{remark}

\subsection*{Entorno a temperatura constante}

Si el entorno es un foco térmico ideal, $S(t)\equiv S_0$, y \eqref{eq:newton} es lineal de coeficientes constantes: $\dot T + kT = kS_0$.

\begin{proposition}[solución del problema de Newton]\label{prop:newton}
La solución del problema \eqref{eq:newton} con $S\equiv S_0$ es
\begin{equation}
  T(t) = S_0 + (T_0 - S_0)\,\ee^{-kt},
  \label{eq:newton-sol}
\end{equation}
y $T(t)\to S_0$ cuando $t\to\infty$.
\end{proposition}

\begin{proof}
La ecuación homogénea $\dot T + kT=0$ tiene solución general $T_h = C\,\ee^{-kt}$, y $T_p\equiv S_0$ es una solución particular. Luego $T = S_0 + C\,\ee^{-kt}$, y la condición inicial $T(0)=T_0$ da $C=T_0-S_0$. Como $k>0$, el término exponencial tiende a cero.
\end{proof}

\begin{remark}
La magnitud $\tau = 1/k$ es la \emph{constante de tiempo}: la diferencia $T-S_0$ se reduce en un factor $\ee^{-1}$ cada $\tau$.
\end{remark}

\begin{figure}[H]
  \centering
  \begin{tikzpicture}
    \begin{axis}[informe, xmin=0, xmax=120, ymin=0, ymax=120,
      xlabel={$t$ (\si{\minute})}, ylabel={$T$ (\si{\celsius})},
      legend style={at={(0.97,0.62)}}]
      \addplot[asintota, forget plot, domain=0:120, samples=2] {20};
      \node[font=\small, text=gris, anchor=south east] at (axis cs:118,21) {$S_0$};
      \addplot[serie1, domain=0:120, samples=100] {20 + 90*exp(-0.04*x)};
      \addlegendentry{$T_0>S_0$}
      \addplot[serie2, domain=0:120, samples=100] {20 - 15*exp(-0.04*x)};
      \addlegendentry{$T_0<S_0$}
    \end{axis}
  \end{tikzpicture}
  \caption{Ley de Newton con $S_0=\SI{20}{\celsius}$ y $k=\SI{0,04}{\per\minute}$. Toda solución converge a $S_0$.}
  \label{fig:newton}
\end{figure}

\begin{example}
Un cuerpo a $T_0=\SI{90}{\celsius}$ se sitúa en un entorno a $S_0=\SI{20}{\celsius}$. Tras \SI{10}{\minute} su temperatura es \SI{60}{\celsius}. Determinar $k$ y el instante en que alcanza \SI{30}{\celsius}.

De \eqref{eq:newton-sol}, $60 = 20 + 70\,\ee^{-10k}$, luego
\[
  k = \tfrac{1}{10}\ln\tfrac{70}{40} \approx \SI{0,056}{\per\minute}.
\]
Imponiendo $T(t)=30$: $\ee^{-kt}=\tfrac{10}{70}$, de donde
$t = \tfrac{1}{k}\ln 7 \approx \SI{34,8}{\minute}$.\finejemplo
\end{example}

\newpage

\section{Dinámica de poblaciones}

Sea $P(t)$ el tamaño de una población. Los modelos de esta sección son ecuaciones autónomas $\dot P = f(P)$, donde $f(P)=P\,g(P)$ y $g$ es la tasa de crecimiento per cápita. Los modelos difieren en la forma de $g$.

\subsection{El modelo de Malthus}

\begin{definicion}[Modelo de Malthus]{Modelo de Malthus}
La tasa de crecimiento per cápita es constante, $g\equiv k$:
\begin{equation}
  \dv{P}{t} = kP, \qquad P(0)=P_0>0,
  \label{eq:malthus}
\end{equation}
donde $k=\text{natalidad}-\text{mortalidad}$.
\end{definicion}

Por separación de variables, $\int \dd P/P = \int k\,\dd t$, y por tanto
\begin{equation}
  P(t) = P_0\,\ee^{kt}.
  \label{eq:malthus-sol}
\end{equation}

Si $k>0$ la población crece exponencialmente, con tiempo de duplicación $t_2=\ln 2/k$; si $k<0$ se extingue exponencialmente. En ambos casos el modelo carece de límite superior, lo que lo hace inválido a largo plazo cuando los recursos son finitos.

\subsection{El modelo logístico}

\begin{definicion}[Modelo logístico]{Modelo logístico (Verhulst)}
La tasa per cápita decrece linealmente con $P$ y se anula en la capacidad de carga $M>0$:
\begin{equation}
  \dv{P}{t} = rP\Bigl(1-\frac{P}{M}\Bigr), \qquad P(0)=P_0>0,
  \label{eq:logistica}
\end{equation}
con $r>0$ la tasa de crecimiento intrínseca.
\end{definicion}

La ecuación \eqref{eq:logistica} es de Bernoulli con exponente $2$. Sus equilibrios son $P=0$ (inestable) y $P=M$ (asintóticamente estable).

\begin{proposition}[solución logística]\label{prop:logistica}
La solución de \eqref{eq:logistica} es
\begin{equation}
  P(t) = \frac{M P_0}{P_0 + (M-P_0)\,\ee^{-rt}},
  \label{eq:logistica-sol}
\end{equation}
y $P(t)\to M$ cuando $t\to\infty$ para todo $P_0>0$.
\end{proposition}

\begin{proof}
Con $z=1/P$ se tiene $\dot z = -\dot P/P^2$, y \eqref{eq:logistica} se convierte en la ecuación lineal
\[
  \dot z + r z = \frac{r}{M}.
\]
Su solución es $z(t)=\tfrac1M + C\,\ee^{-rt}$, y $z(0)=1/P_0$ da $C=\tfrac{M-P_0}{MP_0}$. Entonces $P=1/z$ coincide con \eqref{eq:logistica-sol}.
\end{proof}

\begin{algoritmo}[ecuaciones de Bernoulli $y'+a(x)y=b(x)y^n$]
\begin{pasos}
  \item Reducir a la forma $y'+a\,y=b\,y^n$ con $n\neq 0,1$.
  \item Sustituir $z=y^{1-n}$, de modo que $z'=(1-n)y^{-n}y'$.
  \item Resolver la ecuación lineal $z'+(1-n)a\,z=(1-n)b$.
  \item Deshacer el cambio: $y=z^{1/(1-n)}$, y ajustar la condición inicial.
\end{pasos}
\end{algoritmo}

\begin{figure}[H]
  \centering
  \begin{tikzpicture}
    \begin{axis}[informe, xmin=0, xmax=10, ymin=0, ymax=170,
      xlabel={$t$}, ylabel={$P$},
      legend style={at={(0.97,0.55)}}]
      \addplot[asintota, forget plot, domain=0:10, samples=2] {100};
      \node[font=\small, text=gris, anchor=north west] at (axis cs:0.2,98) {$M$};
      \addplot[serie1, domain=0:10, samples=100] {1000/(10+90*exp(-0.8*x))};
      \addlegendentry{$P_0<M$}
      \addplot[serie2, domain=0:10, samples=100] {16000/(160-60*exp(-0.8*x))};
      \addlegendentry{$P_0>M$}
    \end{axis}
  \end{tikzpicture}
  \caption{Soluciones logísticas con $M=100$, $r=0{,}8$. Si $P_0<M$ la curva es sigmoidea; si $P_0>M$ decrece monótonamente hacia $M$.}
  \label{fig:logistica}
\end{figure}

\subsection{La ley de Gompertz}

\begin{definicion}[Ley de Gompertz]{Ley de Gompertz}
La tasa per cápita es proporcional al logaritmo del cociente entre la capacidad de carga y la población:
\begin{equation}
  \dv{P}{t} = bP\,\ln\frac{M}{P}, \qquad P(0)=P_0>0,
  \label{eq:gompertz}
\end{equation}
con $b>0$ y $M>0$ la capacidad de carga.
\end{definicion}

Con $z=\ln P$ se obtiene $\dot z = b(\ln M - z)$, que es lineal. Su solución es $z=\ln M + \ln(P_0/M)\,\ee^{-bt}$, y por tanto
\begin{equation}
  P(t) = M\exp\!\Bigl[\ln\frac{P_0}{M}\,\ee^{-bt}\Bigr].
  \label{eq:gompertz-sol}
\end{equation}

\begin{lemma}[punto de inflexión]\label{lem:gompertz}
Si $P_0<M/\ee$, la solución \eqref{eq:gompertz-sol} tiene un único punto de inflexión, en $P=M/\ee$.
\end{lemma}

\begin{proof}
Derivando \eqref{eq:gompertz},
\[
  \ddot P = b\dot P\ln\frac{M}{P} - b\dot P = b\,\dot P\Bigl(\ln\frac{M}{P}-1\Bigr).
\]
Como $\dot P>0$ para $0<P<M$, $\ddot P=0$ si y solo si $\ln(M/P)=1$, es decir, $P=M/\ee$.
\end{proof}

\begin{remark}
En la logística la inflexión ocurre en $P=M/2$; en Gompertz, en $P=M/\ee\approx 0{,}37\,M$, de modo que el crecimiento es asimétrico respecto a la sigmoide logística. Este modelo se usa, por ejemplo, en el crecimiento tumoral.
\end{remark}

\begin{figure}[H]
  \centering
  \begin{tikzpicture}
    \begin{axis}[informe, xmin=0, xmax=10, ymin=0, ymax=110,
      xlabel={$t$}, ylabel={$P$},
      legend style={at={(0.97,0.45)}}]
      \addplot[asintota, forget plot, domain=0:10, samples=2] {100};
      \addplot[serie1, domain=0:10, samples=100] {100*exp(ln(0.03)*exp(-0.5*x))};
      \addlegendentry{$P_0/M=0{,}03$}
      \addplot[serie2, domain=0:10, samples=100] {100*exp(ln(0.2)*exp(-0.5*x))};
      \addlegendentry{$P_0/M=0{,}2$}
      \addplot[serie3, domain=0:10, samples=100] {100*exp(ln(0.5)*exp(-0.5*x))};
      \addlegendentry{$P_0/M=0{,}5$}
    \end{axis}
  \end{tikzpicture}
  \caption{Soluciones de Gompertz con $M=100$, $b=0{,}5$.}
  \label{fig:gompertz}
\end{figure}

\subsection{Curvas de infección SIS}

En el modelo SIS, una población de tamaño fijo $N$ se divide en susceptibles $S(t)$ e infectados $I(t)$, y los infectados que se recuperan vuelven a ser susceptibles.

\begin{modelo}[SIS]
\begin{equation}
  \dot S = -\beta SI + \gamma I, \qquad \dot I = \beta SI - \gamma I,
  \label{eq:sis}
\end{equation}
con $\beta>0$ la tasa de contagio y $\gamma>0$ la tasa de recuperación. Se cumple $\dot S+\dot I=0$, luego $S+I=N$.
\end{modelo}

Sustituyendo $S=N-I$ se obtiene una única ecuación de Bernoulli con $n=2$:
\begin{equation}
  \dot I = (\beta N-\gamma)\,I - \beta I^2, \qquad I(0)=I_0\in(0,N].
  \label{eq:sis-I}
\end{equation}
Se define el número reproductivo básico $R_0=\beta N/\gamma$ y $k = N-\gamma/\beta$, de modo que $\beta N-\gamma=\beta k$.

\begin{theorem}[curva de infección SIS]\label{thm:sis}
Para $k\neq 0$, la solución de \eqref{eq:sis-I} es
\begin{equation}
  I(t)=\frac{k}{1+\Bigl(\dfrac{k}{I_0}-1\Bigr)\ee^{-\beta k t}}.
  \label{eq:sis-sol}
\end{equation}
Si $R_0>1$ (es decir, $k>0$), $I(t)\to k$; si $R_0<1$ ($k<0$), $I(t)\to 0$. Si $R_0=1$, $I(t)=I_0/(1+\beta I_0 t)\to 0$.
\end{theorem}

\begin{proof}
Con $z=1/I$, la ecuación \eqref{eq:sis-I} se reduce a $\dot z+\beta k\,z=\beta$. La solución es $z=\tfrac1k+C\,\ee^{-\beta kt}$, y $z(0)=1/I_0$ da $C=\tfrac1{I_0}-\tfrac1k$. Invirtiendo, $I=1/z$ coincide con \eqref{eq:sis-sol}. El comportamiento asintótico se sigue del signo de $k$ en el exponente. Para $k=0$, $\dot z=\beta$ y $z=1/I_0+\beta t$.
\end{proof}

\begin{remark}
El umbral $R_0=1$ separa dos regímenes: la enfermedad persiste con $I_\infty=N(1-1/R_0)$ si $R_0>1$ y se extingue si $R_0\le1$.
\end{remark}

\begin{figure}[H]
  \centering
  \begin{tikzpicture}
    \begin{axis}[informe, xmin=0, xmax=30, ymin=0, ymax=1,
      xlabel={$t$}, ylabel={$I/N$},
      legend style={at={(0.97,0.5)}}]
      % N=1, gamma=0.2: R0=3 (beta=0.6, k=2/3), R0=1.5 (beta=0.3, k=1/3), R0=0.8 (beta=0.16, k=-0.25)
      \addplot[serie1, domain=0:30, samples=100] {(2/3)/(1+((2/3)/0.05-1)*exp(-0.6*(2/3)*x))};
      \addlegendentry{$R_0=3$}
      \addplot[serie2, domain=0:30, samples=100] {(1/3)/(1+((1/3)/0.05-1)*exp(-0.3*(1/3)*x))};
      \addlegendentry{$R_0=1{,}5$}
      \addplot[serie3, domain=0:30, samples=100] {(-0.25)/(1+((-0.25)/0.3-1)*exp(0.16*0.25*x))};
      \addlegendentry{$R_0=0{,}8$}
    \end{axis}
  \end{tikzpicture}
  \caption{Curvas SIS con $\gamma=0{,}2$, $N=1$. Con $R_0>1$ la fracción de infectados tiende a $1-1/R_0$; con $R_0<1$ se extingue.}
  \label{fig:sis}
\end{figure}

\Needspace{12\baselineskip}
\section{Competición de especies}\label{sec:competicion}

\subsection{Modelos de Lotka--Volterra (equilibrio ecológico)}

Consideremos dos especies que interactúan como presa y depredador: los conejos (presa), cuya población es $x(t)$, y los zorros (depredador), con población $y(t)$. Los conejos se reproducen mucho más rápido que los zorros.

\begin{modelo}[Lotka--Volterra]
\begin{equation}
  \dot x = \alpha x - \beta xy, \qquad \dot y = -\gamma y + \delta xy,
  \label{eq:lv}
\end{equation}
con $\alpha,\beta,\gamma,\delta>0$. Los términos $\alpha x$ y $-\gamma y$ son las tasas de reproducción (crecimiento de las presas sin depredadores y decaimiento de los depredadores sin presas), y los términos $\mp xy$ describen la interacción: es mala para los conejos que haya zorros, y buena para los zorros que haya conejos.
\end{modelo}

Dividiendo las dos ecuaciones se elimina el tiempo y se obtiene la ecuación de las órbitas:
\begin{equation}
  \dv{y}{x} = \frac{y\,(\delta x-\gamma)}{x\,(\alpha-\beta y)},
  \label{eq:lv-orbitas}
\end{equation}
que es de variables separables. Separando,
\[
  \frac{\alpha-\beta y}{y}\,\dd y = \frac{\delta x-\gamma}{x}\,\dd x
  \;\Longleftrightarrow\;
  \Bigl(\frac{\alpha}{y}-\beta\Bigr)\dd y + \Bigl(\frac{\gamma}{x}-\delta\Bigr)\dd x = 0 .
\]
La forma de Pfaff resultante es exacta, y se integra directamente:
\begin{equation}
  \alpha\ln y-\beta y + \gamma\ln x-\delta x = \text{cte}
  \qquad\Longleftrightarrow\qquad
  H(x,y)\equiv\delta x-\gamma\ln x+\beta y-\alpha\ln y = c .
  \label{eq:lv-integral}
\end{equation}
La constante se determina con las condiciones iniciales de población: $c=H(x_0,y_0)$.

\begin{theorem}[órbitas cerradas]\label{thm:lv}
El sistema \eqref{eq:lv} tiene dos equilibrios: $(0,0)$, que es un punto de silla, y $(x^\ast,y^\ast)=(\gamma/\delta,\,\alpha/\beta)$. La función $H$ es una integral primera, y sus conjuntos de nivel $\{H=c\}$ en el primer cuadrante, con $c>H(x^\ast,y^\ast)$, son curvas cerradas que rodean a $(x^\ast,y^\ast)$. Por tanto, toda solución con $x_0,y_0>0$ distinta del equilibrio es periódica.
\end{theorem}

\begin{proof}
A lo largo de las soluciones,
\[
  \dot H=\Bigl(\delta-\frac{\gamma}{x}\Bigr)\dot x+\Bigl(\beta-\frac{\alpha}{y}\Bigr)\dot y
  =(\delta x-\gamma)(\alpha-\beta y)+(\beta y-\alpha)(\delta x-\gamma)=0 .
\]
Además, $\nabla H=0$ solo en $(x^\ast,y^\ast)$, donde la matriz hessiana $\operatorname{diag}(\gamma/x^2,\ \alpha/y^2)$ es definida positiva, y $H\to+\infty$ cuando $(x,y)$ se aproxima a la frontera del primer cuadrante o a infinito. Luego $(x^\ast,y^\ast)$ es el único mínimo de $H$ y los conjuntos de nivel son compactos y simples, es decir, curvas cerradas. Como no contienen equilibrios, la solución las recorre periódicamente.
\end{proof}

\begin{desarrollo}[las cuatro fases del ciclo]
Las rectas $x=x^\ast$ e $y=y^\ast$ dividen el plano en cuatro regiones, según el signo de $\dot x=x(\alpha-\beta y)$ y de $\dot y=y(\delta x-\gamma)$. Una órbita se recorre en sentido antihorario:
\begin{itemize}
  \item[\textbf{(I)}] $x>x^\ast$, $y<y^\ast$: $\dot x>0$, $\dot y>0$. Aumentan conejos y zorros.
  \item[\textbf{(II)}] $x>x^\ast$, $y>y^\ast$: $\dot x<0$, $\dot y>0$. Los zorros se comen a los conejos, que disminuyen.
  \item[\textbf{(III)}] $x<x^\ast$, $y>y^\ast$: $\dot x<0$, $\dot y<0$. Hay pocos conejos y los zorros pasan hambre.
  \item[\textbf{(IV)}] $x<x^\ast$, $y<y^\ast$: $\dot x>0$, $\dot y<0$. Como hay menos zorros, los conejos pueden reproducirse mejor.
\end{itemize}
Estos son los llamados \emph{ciclos ecológicos}: ninguna de las especies desaparece.
\end{desarrollo}

\begin{remark}[principio de Volterra]
Si $T$ es el periodo, integrando $\dot y/y=-\gamma+\delta x$ sobre un ciclo se obtiene $0=-\gamma T+\delta\int_0^T x\,\dd t$. Las poblaciones medias son, por tanto, $\langle x\rangle=\gamma/\delta$ y $\langle y\rangle=\alpha/\beta$, es decir, las del equilibrio.
\end{remark}

\begin{figure}[H]
  \centering
  \begin{subfigure}[t]{0.52\linewidth}
    \centering
    \begin{tikzpicture}
      \begin{axis}[informe, width=\linewidth, height=6.4cm, xmin=0, xmax=3.3, ymin=0, ymax=3.3,
        xlabel={$x$ (presas)}, ylabel={$y$ (depredadores)}]
        \addplot[quiver={u=\thisrow{u},v=\thisrow{v},scale arrows=0.2},
                 -{Stealth[length=2.4pt]}, color=azulClaro, forget plot]
          table[x=x,y=y]{figuras/datos/lv_campo.dat};
        \addplot[serie3, forget plot] table[x=x,y=y]{figuras/datos/lv_orbita_c.dat};
        \addplot[serie2, forget plot] table[x=x,y=y]{figuras/datos/lv_orbita_b.dat};
        \addplot[serie1, forget plot] table[x=x,y=y]{figuras/datos/lv_orbita_a.dat};
        \addplot[asintota, forget plot, domain=0:3.3, samples=2] {1};
        \addplot[asintota, forget plot] coordinates {(1,0) (1,3.3)};
        \addplot[only marks, mark=*, mark size=1.8pt, color=azulOscuro, forget plot] coordinates {(1,1)};
        \node[font=\scriptsize, anchor=north west, text=azulOscuro, fill=white, inner sep=1pt] at (axis cs:1.03,0.97) {$(x^\ast,y^\ast)$};
        \node[font=\footnotesize, text=gris] at (axis cs:2.15,0.62) {I};
        \node[font=\footnotesize, text=gris] at (axis cs:2.15,1.75) {II};
        \node[font=\footnotesize, text=gris] at (axis cs:0.5,1.75) {III};
        \node[font=\footnotesize, text=gris] at (axis cs:0.5,0.62) {IV};
      \end{axis}
    \end{tikzpicture}
    \caption{Órbitas $H=c$ ($c=2{,}3;\ 2{,}8;\ 3{,}4$) y campo de direcciones.}
  \end{subfigure}\hfill
  \begin{subfigure}[t]{0.44\linewidth}
    \centering
    \begin{tikzpicture}
      \begin{axis}[informe, width=\linewidth, height=6.4cm, xmin=0, xmax=14, ymin=0, ymax=3,
        xlabel={$t$}, ylabel={Población}, legend columns=1, legend style={at={(0.5,1.04)}, anchor=south}]
        \addplot[serie1] table[x=t,y=x]{figuras/datos/lv_tiempo.dat};
        \addlegendentry{$x$ (presas)}
        \addplot[serie4] table[x=t,y=y]{figuras/datos/lv_tiempo.dat};
        \addlegendentry{$y$ (depredadores)}
      \end{axis}
    \end{tikzpicture}
    \caption{Evolución temporal ($c=2{,}8$): oscilaciones desfasadas.}
  \end{subfigure}
  \caption{Modelo de Lotka--Volterra con $\alpha=\beta=\gamma=\delta=1$. Las órbitas, recorridas en sentido antihorario, son curvas cerradas (ciclos ecológicos) alrededor del equilibrio $(1,1)$.}
  \label{fig:lv}
\end{figure}

\subsection{Curvas de persecución}\label{sec:competicion-pers}

Un conejo se mueve en línea recta y un zorro intenta alcanzarlo, dirigiéndose en cada instante hacia él.

\begin{modelo}[persecución]
Se toma el eje $y$ como la recta a lo largo de la cual se mueve el conejo con velocidad constante $\hat v$, partiendo del origen. Su posición es $\hat{\bm r}=(\hat x,\hat y)$ con
\begin{equation}
  \hat x=0,\qquad \hat y=\hat v t .
  \label{eq:conejo}
\end{equation}
El zorro parte de $(a,0)$ y corre con velocidad $v$, constante en módulo pero no en dirección. Su posición es $\bm r=(x,y)$, con
\begin{equation}
  x(0)=a,\quad y(0)=0,\quad \dot{\bm r}(0)=-v\,\bm\imath .
  \label{eq:zorro-ci}
\end{equation}
En todo instante, la velocidad del zorro apunta hacia el conejo.
\end{modelo}

\begin{figure}[H]
  \centering
  \begin{tikzpicture}
    \begin{axis}[informe, width=0.62\linewidth, height=6.2cm, xmin=-0.08, xmax=1.12, ymin=-0.05, ymax=0.8,
      xlabel={$x$}, ylabel={$y$}, xtick={0,0.5,1}, xticklabels={$0$,$x$,$a$}, ytick=\empty,
      axis lines=left, grid=none]
      \addplot[serie1, domain=0.0001:1, samples=120] {0.5*(x^1.5/1.5 - x^0.5/0.5) + 0.5/0.75};
      \addplot[serie3, forget plot, dashed] coordinates {(0.5,0.0774) (0,0.2542)};
      \addplot[asintota, forget plot] coordinates {(0,0) (0,0.78)};
      \addplot[only marks, mark=*, mark size=1.8pt, color=azulOscuro, forget plot] coordinates {(0.5,0.0774) (1,0)};
      \addplot[only marks, mark=square*, mark size=1.8pt, color=acento, forget plot] coordinates {(0,0.2542)};
      \draw[-{Stealth}, color=acento, thick] (axis cs:0,0.2542) -- (axis cs:0,0.42);
      \node[font=\footnotesize, anchor=west, text=acento] at (axis cs:0.02,0.36) {$\hat v$};
      \node[font=\footnotesize, anchor=east, text=acento] at (axis cs:-0.01,0.2542) {$\hat y=\hat v t$};
      \node[font=\footnotesize, anchor=north east, text=azulOscuro] at (axis cs:0.5,0.065) {$(x,y)$};
      \node[font=\footnotesize, anchor=south, text=azulOscuro] at (axis cs:1,0.02) {zorro};
      \node[font=\footnotesize, anchor=west, text=gris] at (axis cs:0.55,0.24) {línea de visión};
      \draw[-{Stealth}, color=azulOscuro, thick] (axis cs:0.5,0.0774) -- (axis cs:0.40,0.1130);
      \node[font=\footnotesize, anchor=south, text=azulOscuro] at (axis cs:0.36,0.135) {$\bm v$};
    \end{axis}
  \end{tikzpicture}
  \caption{Geometría de la persecución ($k=1/2$). La velocidad del zorro es tangente a su trayectoria y apunta al conejo, de modo que $y'=-(\hat y-y)/x$.}
  \label{fig:persecucion-geo}
\end{figure}

\begin{desarrollo}[ecuación de la trayectoria del zorro]
Sea $y'=\dd y/\dd x$.
\begin{enumerate}[label=\textbf{\arabic*.}, leftmargin=1.6em]
  \item \emph{Módulo de la velocidad.} $v^2=\dot x^2+\dot y^2=\dot x^2\,(1+y'^2)$. Como $x$ decrece con $t$, $\dot x<0$ y
  \begin{equation}
    v\,\dd t=-\sqrt{1+y'^2}\,\dd x .
    \label{eq:pers-modulo}
  \end{equation}
  \item \emph{Dirección.} El zorro corre hacia el conejo: $y'=-\dfrac{\hat y-y}{x}=-\dfrac{\hat vt-y}{x}$, es decir,
  \begin{equation}
    \hat v\,t=y-x\,y' .
    \label{eq:pers-direccion}
  \end{equation}
  \item \emph{Eliminación del tiempo.} Diferenciando \eqref{eq:pers-direccion} respecto a $x$: $\hat v\,\dd t=\bigl(y'-y'-xy''\bigr)\dd x=-x\,y''\,\dd x$. Con \eqref{eq:pers-modulo}, y llamando $k=\hat v/v$,
  \begin{equation}
    x\,y''=k\sqrt{1+y'^2},\qquad k=\frac{\hat v}{v}.
    \label{eq:pers-ode}
  \end{equation}
  \item \emph{Reducción de orden.} Como $y$ no aparece explícitamente, se toma $z=y'$ y queda la ecuación de primer orden separable
  \[
    x\,z'=k\sqrt{1+z^2}\;\Longrightarrow\;\frac{\dd z}{\sqrt{1+z^2}}=k\,\frac{\dd x}{x}.
  \]
\end{enumerate}
\end{desarrollo}

\begin{lemma}[una primitiva elemental]\label{lem:arsinh}
$\displaystyle\int\frac{\dd z}{\sqrt{1+z^2}}=\ln\bigl(z+\sqrt{1+z^2}\,\bigr)+C$.
\end{lemma}

\begin{proof}
Se hace la sustitución trigonométrica $z=\tan\theta$, de modo que $\dd z=\sec^2\theta\,\dd\theta$ y $\sqrt{1+z^2}=\sqrt{1+\tan^2\theta}=\sec\theta$. Entonces
\[
  \int\frac{\sec^2\theta}{\sec\theta}\,\dd\theta=\int\sec\theta\,\dd\theta=\ln\lvert\sec\theta+\tan\theta\rvert+C
  =\ln\bigl(\sqrt{1+z^2}+z\bigr)+C .
\]
(Para integrandos con $\sqrt{a^2-x^2}$, $\sqrt{a^2+x^2}$ o $\sqrt{x^2-a^2}$ se toman, respectivamente, $x=a\sin\theta$, $x=a\tan\theta$ o $x=a\sec\theta$.)
\end{proof}

Integrando, $\ln\bigl(y'+\sqrt{1+y'^2}\bigr)-k\ln x=C$. En $x=a$ se tiene $y'=0$, luego $C=-k\ln a$ y
\begin{equation}
  y'+\sqrt{1+y'^2}=\Bigl(\frac{x}{a}\Bigr)^{k}.
  \label{eq:pers-primera}
\end{equation}
Despejamos $y'$. Aislando la raíz, $\sqrt{1+y'^2}=(x/a)^k-y'$, y elevando al cuadrado, $1=(x/a)^{2k}-2y'(x/a)^k$. De donde
\begin{equation}
  y'=\frac{(x/a)^{2k}-1}{2(x/a)^k}=\frac12\Bigl[\Bigl(\frac{x}{a}\Bigr)^{k}-\Bigl(\frac{x}{a}\Bigr)^{-k}\Bigr].
  \label{eq:pers-yp}
\end{equation}
(Equivalentemente: $\sqrt{1+y'^2}-y'=(x/a)^{-k}$, y se resta o suma a \eqref{eq:pers-primera}.)

\begin{theorem}[trayectoria del zorro]\label{thm:persecucion}
Si $k=\hat v/v<1$, la trayectoria del zorro es
\begin{equation}
  y(x)=\frac a2\left[\frac{1}{1+k}\Bigl(\frac{x}{a}\Bigr)^{1+k}-\frac{1}{1-k}\Bigl(\frac{x}{a}\Bigr)^{1-k}\right]+\frac{ak}{1-k^2},
  \label{eq:pers-sol}
\end{equation}
y el zorro alcanza al conejo en $x=0$, en el instante $\Tf=\dfrac{av}{v^2-\hat v^2}$, en la posición $\hat y_f=y_f=\dfrac{a v\hat v}{v^2-\hat v^2}$.
\end{theorem}

\begin{proof}
Se integra \eqref{eq:pers-yp} con $u=x/a$, $\dd x=a\,\dd u$:
\[
  y=\frac a2\Bigl(\int u^{k}\,\dd u-\int u^{-k}\,\dd u\Bigr)
   =\frac a2\Bigl[\frac{u^{1+k}}{1+k}-\frac{u^{1-k}}{1-k}\Bigr]+C_1 ,
\]
válido para $k\neq1$. La condición $y(a)=0$ ($u=1$) da $C_1=\dfrac a2\Bigl[\dfrac{1}{1-k}-\dfrac{1}{1+k}\Bigr]=\dfrac{ak}{1-k^2}$, lo que prueba \eqref{eq:pers-sol}.

Si $k<1$ el exponente $1-k$ es positivo, de modo que al llegar a $x=0$ los dos primeros términos se anulan y $y(0)=\dfrac{ak}{1-k^2}$. El zorro y el conejo coinciden cuando $x=\hat x=0$ y $y=\hat y$: $\dfrac{ak}{1-k^2}=\hat v\Tf$. Con $k=\hat v/v$,
\[
  \Tf=\frac{ak}{\hat v\,(1-k^2)}=\frac{a}{v}\cdot\frac{1}{1-\hat v^2/v^2}=\frac{av}{v^2-\hat v^2},
  \qquad
  \hat y_f=\hat v\Tf=\frac{av\hat v}{v^2-\hat v^2}.
\]
\end{proof}

\begin{remark}
Si $k>1$ ($\hat v>v$) el término $-(x/a)^{1-k}/(1-k)$ diverge al hacer $x\to0$, y el zorro nunca alcanza al conejo. El caso límite $k=1$ (misma velocidad) se trata en el problema~6 de la \cref{sec:seminarios}: el zorro tampoco lo alcanza y la distancia mínima es $a/2$.
\end{remark}

\begin{figure}[H]
  \centering
  \begin{tikzpicture}
    \begin{axis}[informe, width=0.8\linewidth, height=6.2cm, xmin=0, xmax=1.05, ymin=0, ymax=0.75,
      xlabel={$x/a$}, ylabel={$y/a$}, legend style={at={(0.97,0.92)}, anchor=north east},
      declare function={ pers(\x,\k)=0.5*(\x^(1+\k)/(1+\k)-\x^(1-\k)/(1-\k))+\k/(1-\k*\k); }]
      \addplot[serie1, domain=0.00001:1, samples=150] {pers(x,0.5)};
      \addlegendentry{$v=2\hat v$}
      \addplot[serie2, domain=0.00001:1, samples=150] {pers(x,1/3)};
      \addlegendentry{$v=3\hat v$}
      \addplot[serie4, domain=0.00001:1, samples=150] {pers(x,0.25)};
      \addlegendentry{$v=4\hat v$}
      \addplot[only marks, mark=*, mark size=1.6pt, color=azulOscuro, forget plot]
        coordinates {(0,0.6667) (0,0.375) (0,0.2667)};
    \end{axis}
  \end{tikzpicture}
  \caption{Curvas de persecución del zorro para distintas velocidades relativas, con $a=1$. Los puntos sobre el eje son las posiciones de captura $y_f/a=k/(1-k^2)$.}
  \label{fig:persecucion}
\end{figure}

\begin{example}
Con $a=\SI{100}{\metre}$, $\hat v=\SI{5}{\metre\per\second}$ y $v=\SI{10}{\metre\per\second}$ se tiene $k=1/2$, y
$\Tf=av/(v^2-\hat v^2)=\SI{13,3}{\second}$, con el conejo recorrido $\hat y_f=\hat v\Tf=\SI{66,7}{\metre}$.\finejemplo
\end{example}

\newpage

\Needspace{12\baselineskip}
\section{Fuerzas dependientes de la velocidad}

Una partícula de masa $m$ que se mueve en una dimensión bajo una fuerza que solo depende de la velocidad, $F=F(v)$, queda descrita por la segunda ley de Newton. Ejemplos son el arrastre lineal, $F=-bv$ (bajas velocidades), y el arrastre cuadrático, $F=-Bv^2$ (altas velocidades).

\begin{modelo}[movimiento con $F=F(v)$]
Hay que resolver las dos ecuaciones de primer orden
\begin{equation}
  m\,\dv{v}{t}=F(v),
  \label{eq:vel-newton}
\end{equation}
\begin{equation}
  \dv{x}{t}=v(t).
  \label{eq:vel-posicion}
\end{equation}
Equivalen a la ecuación de segundo orden $m\ddot x=F(\dot x)$, de modo que la solución general $x=\varphi(t,c_1,c_2)$ depende de dos constantes. Estas se determinan con dos condiciones iniciales: $x(t_0)=x_0=\varphi(t_0,c_1,c_2)$ y $\dot x(t_0)=v_0=\varphi'(t_0,c_1,c_2)$.
\end{modelo}

\begin{algoritmo}[movimiento con $F=F(v)$]
\begin{pasos}
  \item \emph{Velocidad.} Separar variables en \eqref{eq:vel-newton}: $\displaystyle\int\frac{\dd v}{F(v)}=\int\frac{\dd t}{m}$. Esto da $t=t(v)$ y, al invertir, $v=v(t)$ (constante $c_1$).
  \item \emph{Posición.} Integrar $\dd x=v(t)\,\dd t$, es decir, $x=x(t)$ (constante $c_2$).
  \item \emph{Condiciones iniciales.} Imponer $v(t_0)=v_0$ y $x(t_0)=x_0$.
  \item \emph{Variante (si interesa $x$ y no $t$).} Con la regla de la cadena, $\dot v=\dv{v}{x}\dv{x}{t}=v\dv{v}{x}$, y \eqref{eq:vel-newton} se escribe $m\,v\,\dd v=F(v)\,\dd x$, ecuación de variables separables que relaciona $x$ y $v$.
\end{pasos}
\end{algoritmo}

\begin{example}[arrastre lineal, $F(v)=-kv$]
Consideramos un movimiento horizontal sin otras fuerzas, con $v(0)=v_0$. Separando variables,
\[
  m\dv{v}{t}=-kv\;\Longrightarrow\;\frac{\dd v}{v}=-\frac{k}{m}\,\dd t .
\]
Si $v(t_f)=v_f<v_0$ (la partícula frena), integrando entre $0$ y $t_f$:
\[
  \int_{v_0}^{v_f}\frac{\dd v}{v}=-\frac km\int_0^{t_f}\dd t
  \;\Longrightarrow\;
  \ln\frac{v_f}{v_0}=-\frac km\,t_f
  \;\Longrightarrow\;
  t_f(v_f)=-\frac mk\ln\frac{v_f}{v_0}.
\]
Como $v_f/v_0<1$, $\ln(v_f/v_0)<0$ y $t_f>0$. Aplicando la exponencial a ambos lados de la ecuación logarítmica, $v_f/v_0=\ee^{-kt_f/m}$, es decir,
\begin{equation}
  v(t)=v_0\,\ee^{-kt/m}.
  \label{eq:vel-lineal}
\end{equation}
La velocidad decrece exponencialmente y tiende a cero, aunque no se anula en tiempo finito. La posición se obtiene integrando \eqref{eq:vel-posicion}:
\[
  x(t)=x_0+\frac{mv_0}{k}\bigl(1-\ee^{-kt/m}\bigr)\;\longrightarrow\;x_0+\frac{mv_0}{k}\quad(t\to\infty),
\]
de modo que la distancia total recorrida es finita, $mv_0/k$.\finejemplo
\end{example}

\begin{example}[arrastre cuadrático, $F(v)=-Bv^2$]
De $m\,\dd v/v^2=-B\,\dd t$ se obtiene $1/v=1/v_0+Bt/m$, es decir,
\[
  v(t)=\frac{v_0}{1+Bv_0t/m},\qquad x(t)=x_0+\frac mB\ln\!\Bigl(1+\frac{Bv_0t}{m}\Bigr).
\]
La velocidad decae como $1/t$ y el recorrido diverge logarítmicamente. Con la variante del algoritmo: $m\,v\,\dd v/\dd x=-Bv^2$ da $\dd v/v=-(B/m)\,\dd x$ y $v=v_0\,\ee^{-B(x-x_0)/m}$.\finejemplo
\end{example}

El caso $F=-k/v^2$, en el que la partícula se detiene en tiempo finito, se resuelve con las dos variantes del algoritmo en el problema~3 de la \cref{sec:seminarios}.

\begin{figure}[H]
  \centering
  \begin{tikzpicture}
    \begin{axis}[informe, xmin=0, xmax=3, ymin=0, ymax=1.05,
      xlabel={$t$ (con $m=k=v_0=1$)}, ylabel={$v/v_0$}, legend style={at={(0.97,0.95)}, anchor=north east}]
      \addplot[serie1, domain=0:3, samples=100] {exp(-x)};
      \addlegendentry{$F=-kv$}
      \addplot[serie2, domain=0:3, samples=100] {1/(1+x)};
      \addlegendentry{$F=-Bv^2$}
      \addplot[serie4, domain=0:0.3333, samples=100] {(1-3*x)^(1/3)};
      \addlegendentry{$F=-k/v^2$}
    \end{axis}
  \end{tikzpicture}
  \caption{Frenado con distintas fuerzas dependientes de la velocidad. Con $F=-k/v^2$ la partícula se detiene en el instante finito $t=mv_0^3/(3k)$.}
  \label{fig:velocidad}
\end{figure}

\newpage

\Needspace{12\baselineskip}
\section{Circuitos eléctricos}

Consideremos un circuito con una resistencia $R$ y una autoinducción $L$ conectadas en serie con una fuente de tensión $U(t)$.

\begin{figure}[H]
  \centering
  \begin{tikzpicture}[american, scale=0.9]
    \draw (0,0) to[sinusoidal voltage source, l=$U(t)$] (0,3)
          to[short] (4,3)
          to[cute inductor, l=$L$, i>^=$I(t)$] (4,0)
          to[R, l=$R$] (0,0);
  \end{tikzpicture}
  \caption{Circuito $RL$ en serie con una fuente de tensión $U(t)$.}
  \label{fig:circuito}
\end{figure}

\begin{modelo}[segunda ley de Kirchhoff]
La intensidad $I(t)$ que circula por el circuito cumple
\begin{equation}
  U(t)=R\,I+L\dv{I}{t},
  \label{eq:kirchhoff}
\end{equation}
donde $U(t)$ es la diferencia de potencial aplicada. Equivale a la ecuación lineal $\dot I+\dfrac RL\,I=\dfrac{U(t)}{L}$, de la forma $\dot I+a(t)I=b(t)$.
\end{modelo}

\begin{remark}[análisis dimensional]
De $[RI]=[L\,\dot I]$ se obtiene $[R]=[L]/[\mathrm T]$, de modo que $[R/L]=[\mathrm T]^{-1}$ tiene dimensiones de frecuencia. La constante de tiempo del circuito es $\tau=L/R$.
\end{remark}

\subsection*{Descarga exponencial de la batería}

Tomemos $U(t)=U_0\,\ee^{-\alpha t}$, con $[\alpha]=[\mathrm T]^{-1}$ (dimensiones de frecuencia). La ecuación a resolver es
\begin{equation}
  R\,I+L\dv{I}{t}=U_0\,\ee^{-\alpha t}
  \;\Longleftrightarrow\;
  \dot I+\frac RL\,I=\frac{U_0}{L}\,\ee^{-\alpha t},
  \qquad I(0)=0 .
  \label{eq:rl-descarga}
\end{equation}

\begin{proposition}[corriente en el circuito $RL$]\label{prop:rl}
Si $\alpha\neq R/L$, la solución de \eqref{eq:rl-descarga} es
\begin{equation}
  I(t)=\frac{U_0}{R-\alpha L}\Bigl(\ee^{-\alpha t}-\ee^{-(R/L)t}\Bigr).
  \label{eq:rl-sol}
\end{equation}
Si $\alpha=R/L$ (\emph{resonancia}), la solución es
\begin{equation}
  I(t)=\frac{U_0}{L}\,t\,\ee^{-(R/L)t}.
  \label{eq:rl-resonancia}
\end{equation}
\end{proposition}

\begin{proof}
\emph{Homogénea.} $\dot I_h+\frac RL I_h=0$ da $\dd I_h/I_h+\frac RL\,\dd t=0$, luego $\ln I_h+\frac RL t=\text{cte}$ e $I_h=c_0\,\ee^{-(R/L)t}$.

\emph{Variación de la constante (Lagrange), $\alpha\neq R/L$.} Se busca $I_p=\kappa(t)\,\ee^{-(R/L)t}$. Al sustituir se cancelan los términos en $\kappa$ y queda
\[
  \dot\kappa\,\ee^{-(R/L)t}=\frac{U_0}{L}\,\ee^{-\alpha t}
  \;\Longrightarrow\;
  \dot\kappa=\frac{U_0}{L}\,\ee^{(R/L-\alpha)t}
  \;\Longrightarrow\;
  \kappa=\frac{U_0}{L}\,\frac{\ee^{(R/L-\alpha)t}}{R/L-\alpha}.
\]
Por tanto, $I_p=\dfrac{U_0}{L(R/L-\alpha)}\,\ee^{-\alpha t}=\dfrac{U_0}{R-\alpha L}\,\ee^{-\alpha t}$, y la solución general es $I=c_0\,\ee^{-(R/L)t}+\dfrac{U_0}{R-\alpha L}\,\ee^{-\alpha t}$. La condición $I(0)=0$ da $c_0=-\dfrac{U_0}{R-\alpha L}$, es decir, \eqref{eq:rl-sol}. La integración solo es posible si $\alpha\neq R/L$.

\emph{Caso $\alpha=R/L$.} Ahora $\dot\kappa=U_0/L$, luego $\kappa=U_0t/L+c_0$, y $I(0)=0$ da $c_0=0$, es decir, \eqref{eq:rl-resonancia}.
\end{proof}

\begin{remark}[resonancia]
En el caso $\alpha=R/L$ la constante de atenuación del estímulo externo coincide con la constante de relajación natural del circuito, $R/L$. Como el término forzante $\ee^{-\alpha t}$ ya es solución de la homogénea, la respuesta se amplifica transitoriamente por un factor lineal $t$, antes de que el decaimiento exponencial acabe por disipar la energía del sistema. La intensidad es máxima en $t=L/R$, con $I_{\mathrm{max}}=U_0/(\ee R)$. Para $\alpha\to R/L$ se puede llegar a \eqref{eq:rl-resonancia} como límite de \eqref{eq:rl-sol}, con la regla de L'Hôpital o con un desarrollo de Taylor (problema~4 de la \cref{sec:seminarios}).
\end{remark}

\begin{remark}[batería constante]
Si $\alpha\to0$ en \eqref{eq:rl-sol} se obtiene $I(t)=\dfrac{U_0}{R}\bigl(1-\ee^{-Rt/L}\bigr)$, que tiende al valor de Ohm, $U_0/R$, con constante de tiempo $\tau=L/R$.
\end{remark}

\begin{figure}[H]
  \centering
  \begin{tikzpicture}
    \begin{axis}[informe, xmin=0, xmax=8, ymin=0, ymax=0.8,
      xlabel={$t$ (con $R/L=1$, $U_0/L=1$)}, ylabel={$I$}, legend columns=2,
      legend style={at={(0.97,0.95)}, anchor=north east},
      declare function={ rl(\t,\a)=(exp(-\a*\t)-exp(-\t))/(1-\a); }]
      \addplot[serie3, domain=0:8, samples=120] {rl(x,0.25)};
      \addlegendentry{$\alpha=\tfrac14\tfrac RL$}
      \addplot[serie2, domain=0:8, samples=120] {rl(x,0.5)};
      \addlegendentry{$\alpha=\tfrac12\tfrac RL$}
      \addplot[serie1, line width=1.5pt, domain=0:8, samples=120] {x*exp(-x)};
      \addlegendentry{$\alpha=\tfrac RL$}
      \addplot[serie4, domain=0:8, samples=120] {rl(x,2)};
      \addlegendentry{$\alpha=2\tfrac RL$}
    \end{axis}
  \end{tikzpicture}
  \caption{Intensidad en el circuito $RL$ con $U=U_0\,\ee^{-\alpha t}$. En resonancia ($\alpha=R/L$) la respuesta es $t\,\ee^{-t}$, con máximo en $t=L/R$.}
  \label{fig:rl}
\end{figure}

\newpage

\Needspace{12\baselineskip}
\section{Seminarios prácticos}\label{sec:seminarios}

Esta sección reúne los problemas del Tema~2 resueltos con detalle y, siempre que es posible, por varios métodos. La \cref{tab:metodos} resume las técnicas empleadas en el tema.

\begin{table}[H]
  \centering
  \small
  \begin{tabular}{@{}p{4.1cm}p{5.6cm}p{3.4cm}@{}}
    \toprule
    \encabezadoTabla{Método} & \encabezadoTabla{Idea} & \encabezadoTabla{Dónde} \\
    \midrule
    Variables separables & $f(y)\,\dd y=g(x)\,\dd x$ & Malthus, Gompertz, \S\ref{sec:poblaciones} \\
    Ecuación lineal: factor integrante & $\ee^{\int a}$ hace exacta $y'+ay=b$ & Newton, circuitos \\
    Variación de la constante (Lagrange) & $y=C(x)\,y_h$ & Newton, circuitos \\
    Coeficientes indeterminados & ensayo $y_p$ de la forma del término forzante; si hay solapamiento con $y_h$, se multiplica por $t$ & Problemas 1 y 2 \\
    Forma de Pfaff y factor integrante & $\mu$ tal que $\mu\omega$ sea exacta & Problema 2, ejemplo final \\
    Riccati / Bernoulli & $P=1/z$ lo linealiza & logística, SIS \\
    Regla de la cadena & $\dot v=v\,\dd v/\dd x$ & Problema 3 \\
    Reducción de orden & $z=y'$ cuando $y$ no aparece & persecución \\
    Límites indeterminados & L'Hôpital, desarrollo de Taylor & Problema 4 \\
    \bottomrule
  \end{tabular}
  \caption{Métodos de resolución utilizados en el Tema~2.}
  \label{tab:metodos}
\end{table}

\subsection{Enunciados}

\begin{problema}
Resolver la ecuación de Newton del enfriamiento en el caso de que la temperatura ambiente evolucione en la forma $S(t)=S_0+b\sin(\omega t)$.
\end{problema}
\begin{problema}
Resolver la ecuación de Newton del enfriamiento en el caso de que la temperatura ambiente evolucione en la forma $S(t)=S_0\,\ee^{-\omega t}$.
\end{problema}
\begin{problema}
Determinar la distancia que recorre hasta pararse un móvil de velocidad inicial $v_0$ sometido a una fuerza dependiente de la velocidad $F=-k/v^2$.
\end{problema}
\begin{problema}
Determinar la evolución temporal de la intensidad en un circuito con resistencia $R$ y autoinducción $L$ con fuerza electromotriz $U=U_0\,\ee^{-\omega t}$ cuando $\omega\to R/L$.
\end{problema}
\begin{problema}
Los elementos radiactivos se desintegran a una tasa proporcional a la cantidad de elemento presente. Proponer un modelo matemático y resolverlo utilizando como parámetros la cantidad inicial de elemento y el periodo de vida media.
\end{problema}
\begin{problema}
Demostrar que un zorro que persigue a un conejo yendo a la misma velocidad que él nunca lo alcanza. Si la distancia inicial entre ambos es $a$, demostrar que la mínima distancia entre ambos es $a/2$.
\end{problema}

\subsection{Soluciones}

\subsubsection*{Problema 1: ambiente oscilante $S(t)=S_0+b\sin\omega t$}

\begin{solucion}
La ley de Newton $\dot T=-k\bigl(T-S(t)\bigr)$ se reordena como la ecuación lineal
\begin{equation}
  \dot T+kT=kS_0+kb\sin\omega t .
  \label{eq:p1}
\end{equation}

\metodo{Método A: coeficientes indeterminados}
\emph{Homogénea.} $\dot T_h+kT_h=0$ da $\dd T_h/T_h+k\,\dd t=0$, $\ln T_h+kt=c$, luego $T_h=C\,\ee^{-kt}$.

\emph{Solución particular.} El término no homogéneo es una constante más una función periódica, así que se propone una solución de la misma naturaleza (sin confundir las constantes con $C$):
\[
  T_p=A+B\sin\omega t+D\cos\omega t,\qquad \dot T_p=B\omega\cos\omega t-D\omega\sin\omega t .
\]
Sustituyendo en \eqref{eq:p1} y agrupando por tipo de función:
\[
  \underbrace{kA}_{\text{cte.}}\;+\;(B\omega+kD)\cos\omega t\;+\;(-D\omega+kB)\sin\omega t
  =kS_0+kb\sin\omega t .
\]
Igualando coeficientes:
\[
  A=S_0,\qquad B\omega+kD=0\;\Rightarrow\;D=-\frac\omega kB,\qquad -D\omega+kB=kb .
\]
Sustituyendo $D$ en la tercera: $B\bigl(\tfrac{\omega^2}{k}+k\bigr)=kb$, es decir, $B\,\tfrac{k^2+\omega^2}{k}=kb$, y por tanto
\[
  B=\frac{k^2b}{k^2+\omega^2},\qquad D=-\frac{k\omega b}{k^2+\omega^2}.
\]
Sumando la homogénea, la solución general es
\begin{equation}
  T(t)=C\,\ee^{-kt}+S_0+\frac{kb}{k^2+\omega^2}\bigl[k\sin\omega t-\omega\cos\omega t\bigr].
  \label{eq:p1-general}
\end{equation}
Con $T(0)=T_0$: $T_0=C+S_0-\dfrac{k\omega b}{k^2+\omega^2}$, de donde $C=T_0-S_0+\dfrac{k\omega b}{k^2+\omega^2}$ y
\begin{equation}
  \boxed{\,T(t)=S_0+\Bigl(T_0-S_0+\frac{k\omega b}{k^2+\omega^2}\Bigr)\ee^{-kt}+\frac{kb}{k^2+\omega^2}\bigl[k\sin\omega t-\omega\cos\omega t\bigr].\,}
  \label{eq:p1-sol}
\end{equation}

\metodo{Método B: factor integrante}
Con la \cref{prop:newton-general}, $T=\ee^{-kt}\bigl[T_0+k\int_0^t\ee^{ks}(S_0+b\sin\omega s)\,\dd s\bigr]$. Usando
\[
  \int \ee^{ks}\sin\omega s\,\dd s=\frac{\ee^{ks}\,(k\sin\omega s-\omega\cos\omega s)}{k^2+\omega^2}
\]
(se obtiene integrando por partes dos veces), se llega de nuevo a \eqref{eq:p1-sol}.

\metodo{Forma amplitud--fase}
Como $k\sin\omega t-\omega\cos\omega t=\sqrt{k^2+\omega^2}\,\sin(\omega t-\varphi)$ con $\tan\varphi=\omega/k$, el régimen estacionario es
\[
  T_p(t)=S_0+\frac{b\,k}{\sqrt{k^2+\omega^2}}\,\sin(\omega t-\varphi):
\]
el objeto oscila con la frecuencia del ambiente, con amplitud atenuada por $k/\sqrt{k^2+\omega^2}<1$ y retardado en fase en $\varphi$. El transitorio $\ee^{-kt}$ decae con constante de tiempo $1/k$.

\metodo{Enfriamiento rápido}
Si $\omega/k\ll1$ (el cuerpo se equilibra mucho más rápido de lo que cambia el ambiente), $\varphi\to0$, la atenuación tiende a $1$ y el coeficiente $k\omega b/(k^2+\omega^2)\to0$. La solución se aproxima por
\[
  T\approx S_0+(T_0-S_0)\,\ee^{-kt}+b\sin\omega t ,
\]
es decir, el cuerpo sigue a la temperatura ambiente $S(t)$ tras un breve transitorio.
\end{solucion}

\begin{figure}[H]
  \centering
  \begin{subfigure}[t]{0.48\linewidth}
    \centering
    \begin{tikzpicture}
      \begin{axis}[informe, width=\linewidth, height=5.4cm, xmin=0, xmax=20, ymin=0, ymax=85,
        xlabel={$t$}, ylabel={$T$}, legend columns=2, legend style={at={(0.5,1.04)}, anchor=south},
        declare function={
          p1T(\t,\k)=20+(80-20+\k*1*8/(\k*\k+1))*exp(-\k*\t)+\k*8/(\k*\k+1)*(\k*sin(deg(\t))-cos(deg(\t)));
        }]
        \addplot[serie3, domain=0:20, samples=300] {20+8*sin(deg(x))};
        \addlegendentry{$S(t)$}
        \addplot[serie1, domain=0:20, samples=300] {p1T(x,10)};
        \addlegendentry{$T(t)$}
      \end{axis}
    \end{tikzpicture}
    \caption{Enfriamiento rápido ($k=10\,\omega$).}
  \end{subfigure}\hfill
  \begin{subfigure}[t]{0.48\linewidth}
    \centering
    \begin{tikzpicture}
      \begin{axis}[informe, width=\linewidth, height=5.4cm, xmin=0, xmax=20, ymin=0, ymax=85,
        xlabel={$t$}, ylabel={$T$}, legend columns=2, legend style={at={(0.5,1.04)}, anchor=south},
        declare function={
          p1T(\t,\k)=20+(80-20+\k*1*8/(\k*\k+1))*exp(-\k*\t)+\k*8/(\k*\k+1)*(\k*sin(deg(\t))-cos(deg(\t)));
        }]
        \addplot[serie3, domain=0:20, samples=300] {20+8*sin(deg(x))};
        \addlegendentry{$S(t)$}
        \addplot[serie1, domain=0:20, samples=300] {p1T(x,0.5)};
        \addlegendentry{$T(t)$}
      \end{axis}
    \end{tikzpicture}
    \caption{Enfriamiento lento ($k=0{,}5\,\omega$): amplitud atenuada y retraso.}
  \end{subfigure}
  \caption{Problema~1. Evolución de $T$ con $S_0=20$, $b=8$, $\omega=1$ y $T_0=80$.}
  \label{fig:p1}
\end{figure}

\subsubsection*{Problema 2: ambiente que decae $S(t)=S_0\,\ee^{-\omega t}$}

\begin{solucion}
La ecuación es $\dot T+kT=kS_0\,\ee^{-\omega t}$. La homogénea es la misma, $T_h=C\,\ee^{-kt}$. El método de coeficientes indeterminados obliga a distinguir dos situaciones, porque el término forzante puede solaparse con la solución de la homogénea cuando $k=\omega$.

\metodo{Caso A: $k\neq\omega$ (sin solapamiento), coeficientes indeterminados}
Se propone $T_p=A\,\ee^{-\omega t}$ (no hace falta incluir una constante: saldría $0$). Entonces $\dot T_p=-\omega A\,\ee^{-\omega t}$ y
\[
  -\omega A+kA=kS_0\;\Longrightarrow\;A=\frac{kS_0}{k-\omega},
\]
luego
\[
  T(t)=C\,\ee^{-kt}+\frac{kS_0}{k-\omega}\,\ee^{-\omega t}.
\]
Con $T(0)=T_0$: $C=T_0-\dfrac{kS_0}{k-\omega}$, y
\begin{equation}
  \boxed{\,T(t)=\frac{kS_0}{k-\omega}\,\ee^{-\omega t}+\Bigl(T_0-\frac{kS_0}{k-\omega}\Bigr)\ee^{-kt}.\,}
  \label{eq:p2-A}
\end{equation}

\metodo{Caso B: $k=\omega$ (resonancia o solapamiento)}
La ecuación es $\dot T+kT=kS_0\,\ee^{-kt}$, y $\ee^{-kt}$ ya forma parte de la solución homogénea. La regla del método dicta multiplicar el ensayo inicial por $t$: $T_p=At\,\ee^{-kt}$, con $\dot T_p=A\,\ee^{-kt}-Akt\,\ee^{-kt}$. Al sustituir se cancelan los términos en $t$ y queda $A\,\ee^{-kt}=kS_0\,\ee^{-kt}$, es decir, $A=kS_0$. La solución general es
\[
  T(t)=C\,\ee^{-kt}+kS_0\,t\,\ee^{-kt}=(C+kS_0t)\,\ee^{-kt},
\]
y con $T(0)=T_0$ resulta $C=T_0$:
\begin{equation}
  \boxed{\,T(t)=(T_0+kS_0t)\,\ee^{-kt}.\,}
  \label{eq:p2-B}
\end{equation}

\metodo{Variación de la constante (Lagrange) en el caso $k=\omega$}
Se parte de $T_h=C\,\ee^{-kt}$ y se hace $C=\hat C(t)$: $T_p=\hat C(t)\,\ee^{-kt}$, $\dot T_p=\dot{\hat C}\,\ee^{-kt}-k\hat C\,\ee^{-kt}$. Al sustituir en la ecuación completa, $\bigl[\dot{\hat C}\,\ee^{-kt}-k\hat C\,\ee^{-kt}\bigr]+k\hat C\,\ee^{-kt}=kS_0\,\ee^{-kt}$. La cancelación de los términos en $\hat C$ es la garantía estructural de que el método va por buen camino. Queda $\dot{\hat C}=kS_0$, $\hat C=kS_0t$ y se recupera $T=(c+kS_0t)\,\ee^{-kt}$.

\metodo{Forma de Pfaff con factor integrante (caso $k\neq\omega$)}
La ecuación $\dot T=-k(T-S_0\,\ee^{-\omega t})$ se escribe como la 1-forma
\[
  \alpha=B\,\dd T+A\,\dd t=1\cdot\dd T+\bigl(kT-kS_0\,\ee^{-\omega t}\bigr)\dd t=0 .
\]
No es exacta: $\partial_tB=0$, pero $\partial_TA=k\neq0$. Se busca un factor integrante $\mu(t)$ tal que $\partial_t(\mu B)=\partial_T(\mu A)$, es decir, $\dot\mu=k\mu$, y se toma $\mu=\ee^{kt}$. La forma
\[
  \dd F=\ee^{kt}\,\dd T+\ee^{kt}\bigl(kT-kS_0\,\ee^{-\omega t}\bigr)\dd t
\]
es exacta, con $F_T=\ee^{kt}$ y $F_t=\ee^{kt}(kT-kS_0\,\ee^{-\omega t})$. Integrando respecto a $t$:
\[
  F(t,T)=\int\ee^{kt}\bigl(kT-kS_0\,\ee^{-\omega t}\bigr)\dd t=T\,\ee^{kt}-\frac{kS_0}{k-\omega}\,\ee^{(k-\omega)t}+\phi(T).
\]
La condición $F_T=\ee^{kt}$ da $\phi'(T)=0$, $\phi=c_1$. Luego $F=\text{cte}$ y, despejando,
\[
  T=\Bigl(\frac{kS_0}{k-\omega}\,\ee^{(k-\omega)t}+C\Bigr)\ee^{-kt},
\]
que coincide con la solución del caso A.

\metodo{Método como ecuación lineal (factor integrante)}
Con \eqref{eq:newton-general}, $T=\ee^{-kt}\bigl[T_0+kS_0\int_0^t\ee^{(k-\omega)s}\,\dd s\bigr]$. El resultado es \eqref{eq:p2-A} si $k\neq\omega$; si $k=\omega$ el integrando es $1$ y resulta \eqref{eq:p2-B}.

\metodo{Comprobación: límite $k\to\omega$}
En \eqref{eq:p2-A}, el término $kS_0\dfrac{\ee^{-\omega t}-\ee^{-kt}}{k-\omega}$ es indeterminado ($0/0$) cuando $k\to\omega$; por L'Hôpital respecto a $k$ tiende a $kS_0\,t\,\ee^{-kt}$, y se recupera \eqref{eq:p2-B}.
\end{solucion}

\begin{figure}[H]
  \centering
  \begin{tikzpicture}
    \begin{axis}[informe, xmin=0, xmax=14, ymin=0, ymax=105,
      xlabel={$t$}, ylabel={$T$}, legend columns=2, legend style={at={(0.97,0.95)}, anchor=north east},
      declare function={ p2T(\t,\k)=(\k*50/(\k-0.5))*exp(-0.5*\t)+(100-\k*50/(\k-0.5))*exp(-\k*\t); }]
      \addplot[serie3, domain=0:14, samples=120] {50*exp(-0.5*x)};
      \addlegendentry{$S(t)$}
      \addplot[serie1, domain=0:14, samples=120] {p2T(x,1)};
      \addlegendentry{$k=2\omega$}
      \addplot[serie2, domain=0:14, samples=120] {(100+0.5*50*x)*exp(-0.5*x)};
      \addlegendentry{$k=\omega$}
      \addplot[serie4, domain=0:14, samples=120] {p2T(x,0.25)};
      \addlegendentry{$k=\omega/2$}
    \end{axis}
  \end{tikzpicture}
  \caption{Problema~2. Respuesta del objeto a un ambiente $S=S_0\,\ee^{-\omega t}$ ($S_0=50$, $\omega=0{,}5$, $T_0=100$). En el caso $k=\omega$ aparece el factor lineal $t$.}
  \label{fig:p2}
\end{figure}

\subsubsection*{Problema 3: distancia de frenado con $F=-k/v^2$}

\begin{solucion}
La segunda ley de Newton da $m\,\dot v=-k/v^2$, con $x(0)=0$ (o $x_0$) y $v(0)=v_0$.

\metodo{Método A: regla de la cadena (la distancia sin pasar por $t$)}
Como se pide la distancia y no el tiempo, se expresa la aceleración en función de la posición:
\[
  a=\dv{v}{t}=\dv{v}{x}\dv{x}{t}=v\dv{v}{x}.
\]
La ecuación se escribe $m\,v\,\dv{v}{x}=-\dfrac{k}{v^2}$, de variables separables: $m\,v^3\,\dd v=-k\,\dd x$. En el estado final el móvil se detiene: $v=0$, $x=d$. Integrando entre las condiciones inicial y final,
\[
  \int_{v_0}^{0}m\,v^3\,\dd v=\int_0^d(-k)\,\dd x
  \;\Longrightarrow\;
  m\Bigl[\frac{v^4}{4}\Bigr]_{v_0}^{0}=-k\,d
  \;\Longrightarrow\;
  -\frac{mv_0^4}{4}=-kd,
\]
de modo que
\begin{equation}
  \boxed{\,d=\frac{m\,v_0^4}{4k}\,}\qquad(\text{más }x_0\text{ si }x(0)=x_0).
  \label{eq:p3-d}
\end{equation}

\metodo{Método B: integrando dos veces}
\emph{Velocidad.} De $v^2\,\dd v=-\dfrac km\,\dd t$ se obtiene $\dfrac{v^3}{3}=-\dfrac km\,t+c$. Con $v(0)=v_0$, $c=v_0^3/3$, luego
\[
  v(t)=\Bigl(v_0^3-\frac{3k}{m}\,t\Bigr)^{1/3}.
\]
\emph{Tiempo de parada.} $v(T)=0$ da $T=\dfrac{mv_0^3}{3k}$.

\emph{Posición.} De $\dot x=\bigl(v_0^3-\tfrac{3k}{m}t\bigr)^{1/3}$,
\[
  x(t)=c'-\frac{m}{4k}\Bigl(v_0^3-\frac{3k}{m}\,t\Bigr)^{4/3},
\]
y con $x(0)=x_0$, $c'=x_0+\dfrac{mv_0^4}{4k}$:
\[
  x(t)=x_0+\frac{mv_0^4}{4k}-\frac{m}{4k}\Bigl(v_0^3-\frac{3k}{m}\,t\Bigr)^{4/3}.
\]
En $t=T$ el paréntesis se anula, y $x(T)=x_0+\dfrac{mv_0^4}{4k}$, que coincide con \eqref{eq:p3-d}.
\end{solucion}

\begin{figure}[H]
  \centering
  \begin{tikzpicture}
    \begin{axis}[informe, xmin=0, xmax=0.4, ymin=0, ymax=1.08,
      xtick={0,0.1,0.2,0.3,0.4}, xlabel={$t/T_\ast$ con $T_\ast=mv_0^3/k$}, ylabel={$v/v_0$, $x/d$}, legend style={at={(0.97,0.55)}, anchor=east}]
      \addplot[asintota, forget plot, domain=0:0.4, samples=2] {1};
      \addplot[serie1, domain=0:0.3333, samples=100] {(1-3*x)^(1/3)};
      \addlegendentry{$v/v_0$}
      \addplot[serie4, domain=0:0.3333, samples=100] {1-(1-3*x)^(4/3)};
      \addlegendentry{$x/d$}
      \addplot[only marks, mark=*, mark size=1.6pt, color=azulOscuro, forget plot] coordinates {(0.3333,0)};
    \end{axis}
  \end{tikzpicture}
  \caption{Problema~3. Velocidad y distancia recorrida (normalizadas); el móvil se detiene en $T=mv_0^3/(3k)$ habiendo recorrido $d=mv_0^4/(4k)$.}
  \label{fig:p3}
\end{figure}

\subsubsection*{Problema 4: circuito $RL$ con $U=U_0\,\ee^{-\omega t}$ cuando $\omega\to R/L$}

\begin{solucion}
Por la segunda ley de Kirchhoff, $RI+L\dot I=U_0\,\ee^{-\omega t}$, es decir, $\dot I+\frac RL I=\frac{U_0}{L}\,\ee^{-\omega t}$, con $I(0)=0$. Para $\omega\neq R/L$, la \cref{prop:rl} da
\[
  I(t)=\frac{U_0}{R-\omega L}\Bigl(\ee^{-\omega t}-\ee^{-(R/L)t}\Bigr),
\]
(homogénea $I_h=c_0\,\ee^{-(R/L)t}$; particular $I_p=\frac{U_0}{R-\omega L}\,\ee^{-\omega t}$ por Lagrange; $I(0)=0$ fija $c_0$). Al hacer $\omega\to R/L$ aparece la indeterminación $0/0$. Se puede resolver de tres formas.

\metodo{Método A: regla de L'Hôpital}
Derivando numerador y denominador respecto a $\omega$:
\[
  \lim_{\omega\to R/L}I(t)=\lim_{\omega\to R/L}\frac{\partial_\omega\bigl[U_0\bigl(\ee^{-\omega t}-\ee^{-(R/L)t}\bigr)\bigr]}{\partial_\omega\,(R-\omega L)}
  =\lim_{\omega\to R/L}\frac{-U_0\,t\,\ee^{-\omega t}}{-L}
  =\frac{U_0}{L}\,t\,\ee^{-(R/L)t}.
\]

\metodo{Método B: desarrollo de Taylor}
Se escribe $I=\dfrac{U_0}{R-\omega L}\,\ee^{-(R/L)t}\Bigl[\ee^{(R/L-\omega)t}-1\Bigr]$. Para $\omega\to R/L$, el desarrollo de primer orden $\ee^{z}\simeq1+z$ con $z=(R/L-\omega)t$ da
\[
  I\simeq\frac{U_0}{R-\omega L}\,\ee^{-(R/L)t}\Bigl(\frac RL-\omega\Bigr)t
  =\frac{U_0}{L}\,t\,\ee^{-(R/L)t},
\]
ya que $R-\omega L=L\,(R/L-\omega)$.

\metodo{Método C: resolución directa del caso resonante}
Con $\omega=R/L$ el término forzante $\ee^{-(R/L)t}$ coincide con la solución de la homogénea. Por Lagrange, $I_p=\kappa(t)\,\ee^{-(R/L)t}$ da $\dot\kappa=U_0/L$, luego $\kappa=U_0t/L$ ($I(0)=0$ implica que la constante es nula). Equivalentemente, el ensayo $I_p=c\,t\,\ee^{-(R/L)t}$ da $c=U_0/L$.

En los tres casos,
\begin{equation}
  \boxed{\,I(t)=\frac{U_0}{L}\,t\,\ee^{-(R/L)t}\,}
  \label{eq:p4}
\end{equation}
\end{solucion}

\begin{figure}[H]
  \centering
  \begin{tikzpicture}
    \begin{axis}[informe, xmin=0, xmax=8, ymin=0, ymax=0.45,
      xlabel={$t$ (con $R/L=1$, $U_0/L=1$)}, ylabel={$I$}, legend columns=2,
      legend style={at={(0.97,0.95)}, anchor=north east},
      declare function={ rl(\t,\a)=(exp(-\a*\t)-exp(-\t))/(1-\a); }]
      \addplot[serie3, domain=0:8, samples=120] {rl(x,0.5)};
      \addlegendentry{$\omega=0{,}5\,R/L$}
      \addplot[serie2, domain=0:8, samples=120] {rl(x,0.8)};
      \addlegendentry{$\omega=0{,}8\,R/L$}
      \addplot[serie4, domain=0:8, samples=120] {rl(x,0.95)};
      \addlegendentry{$\omega=0{,}95\,R/L$}
      \addplot[serie1, line width=1.5pt, domain=0:8, samples=120] {x*exp(-x)};
      \addlegendentry{límite}
    \end{axis}
  \end{tikzpicture}
  \caption{Problema~4. Al aproximar $\omega\to R/L$ las curvas convergen a la respuesta resonante $t\,\ee^{-t}$.}
  \label{fig:p4}
\end{figure}

\subsubsection*{Problema 5: desintegración radiactiva}

\begin{solucion}
Sea $Q(t)$ la cantidad de elemento radiactivo, con $Q(0)=Q_0$. Por hipótesis, la tasa de desintegración es proporcional a la cantidad presente:
\[
  \dv{Q}{t}=-kQ,\qquad k>0 .
\]
Por separación de variables, $Q(t)=Q_0\,\ee^{-kt}$. El periodo de vida media $\tau$ es el tiempo en que la cantidad se reduce a la mitad, $Q(\tau)=Q_0/2$:
\[
  Q_0\,\ee^{-k\tau}=\frac{Q_0}{2}\;\Longrightarrow\;k\tau=\ln2\;\Longrightarrow\;k=\frac{\ln2}{\tau}.
\]
En función de los parámetros pedidos,
\begin{equation}
  \boxed{\,Q(t)=Q_0\,\ee^{-\frac{\ln2}{\tau}\,t}=\frac{Q_0}{2^{\,t/\tau}}\,}
  \label{eq:p5}
\end{equation}
Tras $n$ periodos quedan $Q_0/2^n$. La vida media (tiempo medio de vida de un núcleo) es $1/k=\tau/\ln2\approx1{,}44\,\tau$.
\end{solucion}

\begin{example}
El carbono-14 tiene $\tau\approx\num{5730}$ años. Una muestra conserva el 12,5\,\% de su $^{14}$C inicial cuando $2^{-t/\tau}=2^{-3}$, es decir, $t=3\tau\approx\num{17190}$ años.\finejemplo
\end{example}

\begin{figure}[H]
  \centering
  \begin{tikzpicture}
    \begin{axis}[informe, width=0.8\linewidth, height=5.6cm, xmin=0, xmax=4.2, ymin=0, ymax=1.05,
      xlabel={$t/\tau$}, ylabel={$Q/Q_0$}, xtick={0,1,2,3,4}, ytick={0,0.125,0.25,0.5,1},
      yticklabels={$0$,$\tfrac18$,$\tfrac14$,$\tfrac12$,$1$}]
      \addplot[serie1, domain=0:4.2, samples=100] {2^(-x)};
      \addplot[only marks, mark=*, mark size=1.6pt, color=azulOscuro]
        coordinates {(0,1) (1,0.5) (2,0.25) (3,0.125)};
    \end{axis}
  \end{tikzpicture}
  \caption{Problema~5. Decaimiento radiactivo: cada periodo de vida media la cantidad se reduce a la mitad.}
  \label{fig:p5}
\end{figure}

\subsubsection*{Problema 6: el zorro y el conejo con la misma velocidad}

\begin{solucion}
Con las notaciones de la \cref{sec:competicion-pers}, $k=\hat v/v=1$ y el zorro cumple $x\,y''=\sqrt{1+y'^2}$, con $y(a)=0$ e $y'(a)=0$. Para $k=1$, la solución general \eqref{eq:pers-sol} no es válida (aparecería $1/(1-k)$), pero \eqref{eq:pers-yp} sí lo es:
\[
  p\equiv y'=\frac12\Bigl(\frac xa-\frac ax\Bigr),\qquad
  \sqrt{1+p^2}=\frac12\Bigl(\frac xa+\frac ax\Bigr).
\]
\emph{Trayectoria.} Integrando, $y=\dfrac{x^2}{4a}-\dfrac a2\ln x+C$, y con $y(a)=0$:
\[
  y(x)=\frac{x^2-a^2}{4a}+\frac a2\ln\frac ax .
\]
\emph{Posición del conejo.} Por \eqref{eq:pers-direccion}, $\hat y=\hat vt=y-xy'$. Como $xy'=\dfrac12\Bigl(\dfrac{x^2}{a}-a\Bigr)$,
\[
  \hat y(x)=\frac{a^2-x^2}{4a}+\frac a2\ln\frac ax .
\]
\emph{Distancia entre ambos.} Si $x\to0$, $\hat y\to+\infty$: el conejo se aleja indefinidamente y el instante $t=\hat y/\hat v$ tiende a infinito, sin que el zorro llegue a $x=0$ en tiempo finito. Para cada valor de $x\in(0,a]$, la distancia $d$ cumple
\[
  d^2=(x-\hat x)^2+(y-\hat y)^2=x^2+(xy')^2=x^2+\Bigl(\frac{a^2-x^2}{2a}\Bigr)^{2}=\Bigl(\frac{x^2+a^2}{2a}\Bigr)^{2},
\]
es decir,
\begin{equation}
  d(x)=\frac{x^2+a^2}{2a}>\frac a2\quad\text{para todo }x>0,
  \qquad d\to\frac a2\ \text{cuando }x\to0 .
  \label{eq:p6-d}
\end{equation}
Por tanto, el zorro nunca alcanza al conejo, y la distancia mínima (alcanzada asintóticamente, para $t\to\infty$) es $a/2$.
\end{solucion}

\begin{figure}[H]
  \centering
  \begin{subfigure}[t]{0.48\linewidth}
    \centering
    \begin{tikzpicture}
      \begin{axis}[informe, width=\linewidth, height=5.4cm, xmin=0, xmax=1.05, ymin=0, ymax=2.2,
        xlabel={$x/a$}, ylabel={$y/a,\ \hat y/a$}, legend style={at={(0.97,0.95)}, anchor=north east}]
        \addplot[serie1, domain=0.03:1, samples=150] {(x*x-1)/4+0.5*ln(1/x)};
        \addlegendentry{zorro $y$}
        \addplot[serie4, domain=0.03:1, samples=150] {(1-x*x)/4+0.5*ln(1/x)};
        \addlegendentry{conejo $\hat y$}
      \end{axis}
    \end{tikzpicture}
    \caption{Altura del zorro y del conejo frente a $x$.}
  \end{subfigure}\hfill
  \begin{subfigure}[t]{0.48\linewidth}
    \centering
    \begin{tikzpicture}
      \begin{axis}[informe, width=\linewidth, height=5.4cm, xmin=0, xmax=1.05, ymin=0, ymax=1.05,
        xlabel={$x/a$}, ylabel={$d/a$}]
        \addplot[asintota, forget plot, domain=0:1.05, samples=2] {0.5};
        \addplot[serie1, domain=0:1, samples=100] {(x*x+1)/2};
        \node[font=\small, text=gris, anchor=north west] at (axis cs:0.05,0.48) {$a/2$};
      \end{axis}
    \end{tikzpicture}
    \caption{Distancia $d(x)=(x^2+a^2)/2a$.}
  \end{subfigure}
  \caption{Problema~6 ($k=1$). La distancia entre zorro y conejo decrece hasta $a/2$ sin llegar a anularse.}
  \label{fig:p6}
\end{figure}

\subsection{Complemento: ecuación lineal frente a forma de Pfaff}

Sobre una misma ecuación, estos dos caminos conducen al mismo resultado. Se ilustra con el siguiente ejemplo.

\begin{example}
Resolver $y'-\dfrac{xy}{x^2+1}=\dfrac{x^2-x+1}{x^2+1}\,\ee^x$ como ecuación lineal y buscando un factor integrante.

\metodo{Como ecuación lineal}
\emph{Homogénea.} $y_h'=\dfrac{xy_h}{x^2+1}$ da $y_h=c_0\,(x^2+1)^{1/2}$. \emph{Variación de la constante.} Con $y_p=c(x)\,(x^2+1)^{1/2}$,
\[
  c'=(x^2+1)^{-1/2}\,\ee^x\Bigl(1-\frac{x}{x^2+1}\Bigr)=\dv{}{x}\Bigl[\ee^x\,(x^2+1)^{-1/2}\Bigr]
  \;\Longrightarrow\;c=\frac{\ee^x}{(x^2+1)^{1/2}}\;\Longrightarrow\;y_p=\ee^x .
\]
La solución general es $y=c_0\,(x^2+1)^{1/2}+\ee^x$.

\metodo{Como forma de Pfaff}
La ecuación equivale a $\alpha=\Bigl(\dfrac{xy}{x^2+1}+\dfrac{x^2-x+1}{x^2+1}\,\ee^x\Bigr)\dd x-\dd y=0$. No es exacta. Se busca $\mu(x)$ con $\partial_y(\mu M)=\partial_x(\mu N)$, siendo $M=\dfrac{xy}{x^2+1}+\dfrac{x^2-x+1}{x^2+1}\,\ee^x$ y $N=-1$:
\[
  \mu\,\frac{x}{x^2+1}=-\mu'\;\Longrightarrow\;\mu=(1+x^2)^{-1/2}.
\]
La forma $\dd F=\mu M\,\dd x-\mu\,\dd y$ es exacta, y $F_y=-\mu$ da $F=-(1+x^2)^{-1/2}\,y+\phi(x)$. Imponiendo $F_x=\mu M$ se obtiene $\phi(x)=\ee^x\,(1+x^2)^{-1/2}$. De $F=\text{cte}$:
\[
  y=\ee^x-k\,(1+x^2)^{1/2}.
\]
\finejemplo
\end{example}


% Los ejercicios propuestos se añadirán aquí:
% \input{secciones/07_ejercicios}

\end{document}
